\section{Balanced finite paths near critical frequencies}
\label{sec:balanced-stability}

This section is independent of Gaussian arithmetic. It isolates the
mechanism that turns a finite graph with nonnegative costs into a
quantitative statement about the words of small cost it accepts.

\subsection{One slope controls every factor}

A finite binary word is \emph{balanced} if any two of its factors of the
same length have heavy-letter counts differing by at most one. We denote
the heavy letter by $b$ in this section. The following classical form of
finite balance is useful because the same slope works for every factor.

\begin{lemma}[A compatible slope]
\label{lem:compatible-slope}
For every nonempty finite balanced word $w$, there is an irrational
$\beta\in(0,1)$ such that every factor $v$ satisfies
\begin{equation}
 \bigl||v|_b-\beta|v|\bigr|<1.
 \label{eq:factor-discrepancy}
\end{equation}
\end{lemma}

\begin{proof}
For two nonempty factors $u,v$, put $p=|u|$, $q=|v|$, $A=|u|_b$ and
$B=|v|_b$. We first prove
\begin{equation}
 |qA-pB|<p+q
 \label{eq:balance-determinant}
\end{equation}
by induction on $p+q$. If $p=q$, balance bounds the left side by $p$.
Otherwise interchange the factors so that $p<q$, and write $v=v_1v_2$
with $|v_1|=p$. Then
$$
 qA-pB=p(A-|v_1|_b)+(q-p)A-p|v_2|_b.
$$
Balance bounds the first term in absolute value by $p$, and the induction
hypothesis bounds the remaining two terms together strictly by $q$.
This proves \eqref{eq:balance-determinant}, which after division by $pq$
is the factorial-balance inequality of the classical theory of Sturmian
words \cite[Proposition 2.1.7]{Lothaire2002}.

For each nonempty factor $v$, take the open interval with endpoints
$(|v|_b-1)/|v|$ and $(|v|_b+1)/|v|$. Inequality
\eqref{eq:balance-determinant} says that every two intervals overlap.
There are finitely many intervals, so their largest lower endpoint is
strictly smaller than their smallest upper endpoint. Every lower endpoint
is below one and every upper endpoint is positive. The common intersection
with $(0,1)$ is therefore nonempty and open. Any irrational $\beta$ in
this intersection proves \eqref{eq:factor-discrepancy}.
\end{proof}

The passage from that inequality to a single compatible slope is also
classical; see \cite[Chapter~2]{Lothaire2002}. We included the proof to
keep the finite-path argument self-contained. The compatible slope need not equal the empirical density
$|w|_b/|w|$. For instance, $00100010000$ is balanced, but its factor $10001$
has discrepancy $12/11$ relative to that density. The common slope in
\cref{lem:compatible-slope} avoids this finite-word endpoint error.

\subsection{The slope interval as a projection}

The interval constructed in the proof of \cref{lem:compatible-slope} is not
an artefact of that proof. It is the shadow, on the slope axis, of a planar
region attached to $w$, and the discrepancy condition is exactly what
survives when the second parameter is eliminated.

Read $b$ as $1$ and $a$ as $0$, and recall that the \emph{mechanical word}
with slope $\beta$ and intercept $\gamma$ has $i$th letter
$\lfloor\beta i+\gamma\rfloor-\lfloor\beta(i-1)+\gamma\rfloor$.

\begin{proposition}[The compatible slopes are a projection]
\label{prop:slope-projection}
Let $w=w_1\cdots w_n$ be a nonempty binary word, put
$k_i=\lvert w_1\cdots w_i\rvert_b$ for $0\leqslant i\leqslant n$, and set
\begin{equation}
 \mc F(w)=\bigl\{(\beta,\gamma)\in\R^2:\;
   k_i\leqslant\beta i+\gamma<k_i+1,\quad 0\leqslant i\leqslant n\bigr\}.
 \label{eq:parameter-region}
\end{equation}
Then $\mc F(w)$ is the set of parameters whose mechanical word of length $n$
is $w$, and its image under $(\beta,\gamma)\mapsto\beta$ is
\begin{equation}
 J(w)=\bigl\{\beta\in\R:\;
   \bigl\lvert\,\lvert v\rvert_b-\beta\lvert v\rvert\,\bigr\rvert<1
   \text{ for every nonempty factor }v\text{ of }w\bigr\},
 \label{eq:slope-interval}
\end{equation}
an open interval with rational endpoints. It is nonempty precisely when $w$
is balanced, and the slope of \cref{lem:compatible-slope} is any irrational
point of $J(w)\cap(0,1)$.
\end{proposition}

\begin{proof}
The constraint at $i=0$ reads $0\leqslant\gamma<1$, so on $\mc F(w)$ one has
$\lfloor\beta i+\gamma\rfloor=k_i$ for every $i$, and the successive
differences of the $k_i$ are the letters of $w$; conversely those
differences determine the $k_i$ from $k_0=0$. This is the first assertion.

Eliminate $\gamma$. The system \eqref{eq:parameter-region} has a solution
$\gamma$ for a given $\beta$ exactly when
$$
 \max_{0\leqslant i\leqslant n}\bigl(k_i-\beta i\bigr)
 <\min_{0\leqslant j\leqslant n}\bigl(k_j+1-\beta j\bigr),
$$
that is, when $(k_i-k_j)-\beta(i-j)<1$ for all $i\neq j$, the case $i=j$
being vacuous. For $i>j$ the factor $v=w_{j+1}\cdots w_i$ has
$\lvert v\rvert_b=k_i-k_j$ and $\lvert v\rvert=i-j$, so the condition is
$\lvert v\rvert_b-\beta\lvert v\rvert<1$; for $i<j$ the same factor, taken
with the roles reversed, gives $\beta\lvert v\rvert-\lvert v\rvert_b<1$.
The two families together are \eqref{eq:slope-interval}. Being a finite
intersection of the open intervals with endpoints
$(\lvert v\rvert_b\mp1)/\lvert v\rvert$, the set $J(w)$ is an open interval
with rational endpoints.

If $\beta\in J(w)$ and $u,v$ are factors of a common length $p$, then
$\bigl\lvert\lvert u\rvert_b-\lvert v\rvert_b\bigr\rvert<2$, hence at most
one by integrality, so $w$ is balanced; the converse is
\cref{lem:compatible-slope}, whose proof exhibits a point of $J(w)$ in
$(0,1)$. An open interval contains an irrational point.
\end{proof}

\begin{remark}[The two-parameter picture]
\label{rem:duality}
The regions $\mc F(w)$, for $w$ ranging over the balanced words of length
$n$, are the faces of an arrangement of lines in the $(\beta,\gamma)$ plane:
each constraint in \eqref{eq:parameter-region} is a half-plane, and the
bounding lines are dual to the lattice points involved. Berstel and
Pocchiola used exactly this duality, together with Euler's relation, to
count the faces and so to recover Mignosi's enumeration of the finite
factors of Sturmian words, of which there are
$1+\sum_{i=1}^{n}(n-i+1)\varphi(i)$ of length $n$, with $\varphi$ the Euler
function \cite{BerstelPocchiola1993,Mignosi1991}. In that reading,
\cref{prop:slope-projection} says that
\cref{lem:compatible-slope} computes the shadow of one face, and that the
intercept is the coordinate being integrated out. The two parameters play
complementary roles: over a fixed irrational slope the fibre of the
projection is cut into $n+1$ cells, one for each factor of that length of a
Sturmian word of that slope, whereas moving the slope changes the face, and
so the word itself.
\end{remark}

\subsection{Endpoint weights and critical components}

Let $G$ be a finite directed graph. Each edge reads exactly one letter in
$\{a,b\}$; parallel edges and nondeterminism are allowed. Initial states
have weights $I(s)$, terminal states have weights $T(s)$, and edges have
weights $c(e)$. Assume all these weights are nonnegative integers.
The completed weight of an accepted path is
\begin{equation}
 \Delta=I(s_0)+\sum_{e\text{ on the path}}c(e)+T(s_n).
 \label{eq:abstract-cost}
\end{equation}

These weights may result from a potential transformation. If original
weights are $w_0,w(e),f$ and a potential is $\lambda$, the relevant transformation is
\begin{equation}
 I(s_0)=w_0-\lambda(s_0),\quad
 c(e)=w(e)+\lambda(s)-\lambda(t),\quad
 T(s_n)=\lambda(s_n)+f(s_n).
 \label{eq:weight-pushing}
\end{equation}
This is the usual weight-pushing transformation \cite[Section 6.3]{Mohri2009}. The terms telescope, preserving the completed cost. Nonnegative cycle
weights alone do not guarantee nonnegative values of both transformed
endpoint weights.

Let $G_0$ retain the edges of weight zero. Its strongly connected components
form an acyclic condensation graph $\mathcal D$. Call a component
\emph{cyclic} if it contains a directed cycle, including a self-loop.
Suppose each cyclic component $C_i$ admits a reduced rational number
$\theta_i=p_i/d_i$, with $d_i>0$, and an integer function $H_i$ on its vertices such that
\begin{equation}
 d_i\mathbf1_{\{e\text{ reads }b\}}-p_i=H_i(t)-H_i(s)
 \label{eq:count-potential}
\end{equation}
for every internal edge $e:s\longrightarrow t$. Put
\begin{equation}
 D_i=\frac{\max H_i-\min H_i}{d_i},\qquad
 \Theta=\{\theta_i:C_i\text{ cyclic}\}.
 \label{eq:critical-data}
\end{equation}
Summation around a cycle gives $0\leqslant\theta_i\leqslant1$.
Summation along an internal path gives
\begin{equation}
 \bigl||v|_b-\theta_i|v|\bigr|\leqslant D_i.
 \label{eq:internal-count-bound}
\end{equation}

Condition \eqref{eq:count-potential} is equivalent to all directed cycles
in the component having the same letter frequency. To see the converse,
assign to a vertex the sum of $d_i\mathbf1_b-p_i$ along a path from a fixed
root. Two choices give the same sum: append a common return path, then
decompose the resulting closed walks into cycles. This also shows how to
check the condition by a finite list of edge identities.

Give each cyclic vertex of $\mathcal D$ weight $D_i+1$, every other vertex
weight zero and every condensation edge weight one. Define
\begin{equation}
 L=\max_{\gamma\text{ a directed path in }\mathcal D}
 \left(\#\text{edges of }\gamma+
       \sum_{C_i\in\gamma\text{ cyclic}}(D_i+1)\right),\qquad K=L+1.
 \label{eq:condensation-budget}
\end{equation}
Paths consisting of a single vertex are included. The number $L$ is finite
and is obtained by the usual longest-path recursion on an acyclic graph.

It is convenient to allow a bounded number of letters to be read outside
the graph. Fix an integer $h\geqslant0$, and suppose that an initial
prefix and a final suffix of total length at most $h$ are consumed before
entering and after leaving $G$, while \eqref{eq:abstract-cost} still
records the completed cost of the whole word. The case in which no letter
is omitted is $h=0$.

\subsection{The stability theorem}

We isolate first the counting step, because it uses only one property of a
real number $\beta$, namely that it be compatible with $w$ in the sense of
\eqref{eq:factor-discrepancy}, and not the particular slope furnished by
\cref{lem:compatible-slope}. The arithmetic application in
\cref{sec:markoff-application} will supply its own compatible slope.

\begin{lemma}[The budget inequality]
\label{lem:budget}
Assume $\Theta\ne\varnothing$. Let $w$ be a nonempty word of length $n$, all
but at most $h$ of whose letters are read by an accepted path, the
remaining ones forming a prefix and a suffix consumed outside the graph.
Let $\beta$ be a real number with
$\bigl||v|_b-\beta|v|\bigr|<1$ for every nonempty factor $v$ of $w$, and put
$\delta=\operatorname{dist}(\beta,\Theta)$. If $\delta\leqslant1$, then
\begin{equation}
 n\delta\leqslant h+\Delta+(\Delta+1)L.
 \label{eq:budget}
\end{equation}
\end{lemma}

\begin{proof}
Let $k$ be the number of positive-weight edges in the path. Integrality and
nonnegative endpoint weights imply $k\leqslant\Delta$.

Split the path into $k$ positive edges and $k+1$ possibly empty zero-edge
segments. A zero segment visits each condensation vertex at most once.
Inside a cyclic component, its internal subpath reads a factor $v$ of the
full word, and the hypothesis on $\beta$ together with
\eqref{eq:internal-count-bound} implies
$$
 \delta|v|\leqslant|\beta-\theta_i||v|<1+D_i.
$$
A noncyclic component is a singleton without an internal edge. Every
remaining letter crosses a condensation edge and contributes at most one
to $\delta$ times length, since $\delta\leqslant1$. Thus each zero segment
contributes at most $L$. The positive edges contribute at most $k$, and
the omitted boundary letters at most $h$. Consequently
$n\delta\leqslant h+k+(k+1)L$, which is \eqref{eq:budget} because
$k\leqslant\Delta$ and $L\geqslant0$.
\end{proof}

\begin{theorem}[Balanced finite-path stability]
\label{thm:balanced-stability}
Assume $\Theta\ne\varnothing$. Let $w$ be a balanced word all but at most
$h$ of whose letters are read by an accepted path, the remaining ones
forming a prefix and a suffix consumed outside the graph. Then
\begin{equation}
 \min_{\theta\in\Theta}\bigl||w|_b-\theta|w|\bigr|
 <K(\Delta+1)+h.
 \label{eq:stability-main}
\end{equation}
Both counts on the left refer to the whole word. In particular, when no
letter is omitted the bound is $K(\Delta+1)$.
\end{theorem}

\begin{proof}
The empty word is immediate. Otherwise put $n=|w|$, $m=|w|_b$, and choose
$\beta$ from \cref{lem:compatible-slope}; it satisfies
\eqref{eq:factor-discrepancy}, and
$\delta=\operatorname{dist}(\beta,\Theta)\leqslant1$ because $\beta$ and
every member of $\Theta$ lie in $[0,1]$. \Cref{lem:budget} applies. Choosing
a member of $\Theta$ nearest to $\beta$ and applying
\eqref{eq:factor-discrepancy} also to the whole word gives
$$
 \min_{\theta\in\Theta}|m-\theta n|
 <1+n\delta\leqslant h+(\Delta+1)(L+1)
 \leqslant h+K(\Delta+1).
$$
\end{proof}

\begin{corollary}[Linear cost away from the critical set]
\label{cor:linear-cost}
Under the hypotheses of \cref{thm:balanced-stability}, if $n>0$ and
$\operatorname{dist}(m/n,\Theta)\geqslant\eta>0$, then
$$
 \Delta>\frac{\eta n-h}{K}-1.
$$
Consequently, for a fixed $h$, every accumulation point of the empirical
frequencies of balanced accepted words of bounded completed cost and
unbounded length lies in $\Theta$.
\end{corollary}

If $G_0$ has no cyclic component, one must not take a minimum over an empty
set. In this case each zero segment has at most $c-1$ edges, where $c$ is
the number of components. The corresponding conclusion is the direct
length bound $n\leqslant h+c(\Delta+1)-1$.

\subsection{Why the finite set is stronger than its convex hull}

Assume again that $\Theta\ne\varnothing$. Without balance, let
$$
 L_0=\max_{\gamma\subset\mathcal D}
 \left(\#\text{edges of }\gamma+\sum_{C_i\in\gamma\text{ cyclic}}D_i\right).
$$
Assign frequency $\theta_i$ to every internal cyclic-component segment,
and any fixed member of $\Theta$ to the remaining letters. The
length-weighted average of these assigned frequencies belongs to
$\operatorname{conv}\Theta$. Telescoping the internal count potentials
therefore gives, for an arbitrary nonempty input,
\begin{equation}
 n\operatorname{dist}(m/n,\operatorname{conv}\Theta)
 \leqslant h+k+(k+1)L_0.
 \label{eq:convex-hull-bound}
\end{equation}
If all critical frequencies coincide, this already controls the count
without balance.

In general the finite set cannot replace its convex hull. Take a zero-cost
$a$ loop at an initial state, a zero-cost $b$ loop at a terminal state,
and a zero-cost $a$ edge between them. Both endpoint weights are zero.
The critical set is $\{0,1\}$, yet the accepted words $a^{r+1}b^r$ have
cost zero and count distance $r$ from that set. For $r\geqslant2$ these
words are unbalanced. The example identifies exactly the obstruction
removed by \cref{lem:compatible-slope}.
